\chapter{Hematemesis (Vomiting Blood)}

\section*{Definition}
Hematemesis is the vomiting of blood originating from the upper gastrointestinal tract (proximal to the ligament of Treitz). It may present as bright red blood or as coffee-ground material (partially digested blood). Hematemesis is a medical emergency requiring immediate evaluation.

\section*{What Happens in Your Body}
Upper GI bleeding results from disruption of mucosal integrity exposing underlying blood vessels. Bleeding source determines color and character: arterial bleeding produces bright red blood; venous bleeding or slow hemorrhage produces darker blood that turns coffee-ground after exposure to gastric acid. Significant blood loss activates compensatory mechanisms (tachycardia, narrowing of blood vessels) and can progress to hemorrhagic shock.

\section*{Causes}
\begin{itemize}
  \item Peptic ulcer disease (gastric or duodenal ulcers -- most common cause)
  \item Esophageal varices (portal hypertension in liver cirrhosis)
  \item Mallory-Weiss tear (mucosal tear at gastroesophageal junction from retching/vomiting)
  \item Gastritis or erosive esophagitis (NSAIDs, alcohol, stress-related)
  \item Esophagitis (reflux, infectious, pill-induced)
  \item Gastric or esophageal cancer
  \item Dieulafoy lesion (aberrant submucosal artery)
  \item Angiodysplasia (vascular malformations)
  \item Coagulopathy (anticoagulants, liver disease, hemophilia)
  \item Hemobilia (bleeding from biliary tree)
\end{itemize}

\section*{When to See a Doctor}
See your doctor if:
\begin{itemize}
  \item Any episode of hematemesis requires immediate medical evaluation.
  \item Signs of hemodynamic instability: lightheadedness, syncope, pallor, tachycardia, hypotension
  \item Large-volume hematemesis (massive hemorrhage)
  \item Hematochezia or melena accompanying hematemesis (indicates significant bleeding)
  \item Known liver disease or esophageal varices
  \item Anticoagulant or antiplatelet therapy
  \item Elderly people or those with comorbidities (higher risk of decompensation)
\end{itemize}

\section*{Related Symptoms}
\begin{itemize}
  \item Melena (black, tarry stool)
  \item Hematochezia (bright red blood per rectum -- massive upper GI bleed)
  \item Abdominal pain
  \item Epigastric tenderness
  \item Nausea
  \item Syncope or presyncope
  \item Weakness
  \item Pallor
  \item Tachycardia
  \item Hypotension
  \item Dyspnea
\end{itemize}

\section*{Self-Care / Home Management}
\begin{itemize}
  \item Call emergency services immediately for any hematemesis.
  \item Do not eat or drink until evaluated by a physician.
  \item Lie on the side to prevent aspiration if vomiting.
  \item Collect vomitus in a container for medical inspection.
  \item Stop any aspirin, NSAIDs, or anticoagulants (consult physician).
  \item Do not attempt self-treatment at home.
\end{itemize}

\section*{How Your Doctor Will Evaluate You}
\begin{itemize}
  \item ABCs (airway, breathing, circulation) with intravenous access and fluid resuscitation.
  \item Complete blood count, coagulation profile, liver function tests, type and crossmatch.
  \item Upper endoscopy (EGD) within 24 hours for diagnosis and therapeutic intervention.
  \item Nasogastric tube lavage to assess bleeding activity.
  \item Imaging (CT angiography) if endoscopy nondiagnostic or patient unstable.
  \item Blatchford or Rockall score for risk stratification.
\end{itemize}

\section*{Treatment Options}
\begin{itemize}
  \item Endoscopic therapy: Injection, cautery, clipping, band ligation of bleeding site.
  \item Intravenous PPIs (high-dose bolus followed by infusion) after endoscopic hemostasis.
  \item Octreotide for variceal bleeding.
  \item Transfusion of packed red blood cells for significant anemia (target Hb > 7-8 g/dL).
  \item Surgical intervention or interventional radiology (embolization) for refractory bleeding.
  \item Treat underlying cause: H. pylori eradication, PPI for ulcer, beta-blockers for variceal prophylaxis.
\end{itemize}

\clearpage
